\documentclass[10pt, a4paper]{article}

% --- Packages ---
\usepackage[top=2cm, bottom=2cm, left=2cm, right=2cm]{geometry}
\usepackage{graphicx} % For images
\usepackage{amsmath}  % For math/equations
\usepackage{hyperref} % For clickable links
\usepackage{caption}  % For styling captions
\usepackage{titlesec} % For controlling section title sizes
\usepackage{float}
\usepackage{xcolor} % Required for defining and using colors
\usepackage{hyperref} % Keep this as your last \usepackage
\usepackage{color}
\usepackage{soul}

% ברירת המחדל של המרקור היא צהוב. אם תרצי לשנות למשל לטורקיז/תכלת, הוסיפי גם את השורה הזו:
% \sethlcolor{cyan}

\hypersetup{
    colorlinks=true, % Removes the default boxes and colors the text instead
    linkcolor=blue,  % Color for internal links (e.g., Figure and Table references)
    citecolor=blue,  % Color for bibliographical citations
    urlcolor=blue    % Color for external web URLs
}
% --- Document Formatting ---
\setlength{\parindent}{0pt}
\setlength{\parskip}{0.5em}
\captionsetup{font=scriptsize, labelfont=bf}

% --- Customize Title Sizes ---
\titleformat{\section}
  {\large\bfseries}{\thesection}{1em}{}
\titleformat{\subsection}
  {\normalsize\bfseries}{\thesubsection}{1em}{}

\usepackage{fancyhdr}
\pagestyle{fancy}
\fancyhf{} % מנקה את הגדרות ברירת המחדל
\fancyhead[L]{\footnotesize Weekly Summary \#3 \today} % יופיע למעלה בשמאל
\fancyhead[R]{\footnotesize Rachel Broner $\vert$ Miri Adler's Lab} % יופיע למעלה בימין
\fancyfoot[C]{\thepage} % מספר עמוד למטה באמצע

% \usepackage{mathptmx}

\usepackage{setspace}
\onehalfspacing % For 1.5 line spacing


\begin{document}

\begin{center}
    {\Large \textbf{Weekly Summary \#6 $\vert$ Association Rule Mining}} \\[0.5em]
    {\large Part 2 $\vert$ FOVs as vectors in rule space} \\[0.5em]
    \textit{GVHD Analysis} $\vert$ Week of: \today
\end{center}
\vspace{0.2cm}
\hrule
\vspace{0.5cm}

% ---------------------------------------------------------
\section{FOVs as vectors in rule space}

\textbf{How the rule space is built.}
Every FOV is a row and every rule a column, holding that rule's lift in that FOV; a rule that
never fired is set to $1$, meaning no association. Only rules seen in more than 2\% of the
FOVs become columns. The values are log-transformed and standardised, then the PCA is run.

\textbf{Why 2\% and not 0.} Keeping every rule would leave more rules than FOVs, which makes
a PCA meaningless. 2\% is the widest cut that avoids that, and it only drops rules seen in
one or two FOVs.

Each part below is one PCA, shown the same way: how much each component carries, where the
FOVs fall, which rules pull each component, what its ends look like in tissue, and whether it
merely follows how common a cell type is. A loading is a rule's weight in a component --
green pulls one way, salmon the other. For the cell types, \emph{Spearman's rho} compares the
two orderings of the FOVs: $+1$ the same order, $0$ no relation, $-1$ the reverse; the grey
band marks values too small to reach $p<0.05$. The FOV maps are wide and meant to be zoomed
into.

% ---------------------------------------------------------
\subsection{All FOVs}

Every stage and both organs. The two scatters are the same PCA
coloured twice, so they can be compared directly.

\begin{figure}[H]
    \centering
    \includegraphics[width=0.62\textwidth]{summary_downloads/rs_all_scree.pdf}
    \caption{How much of the picture each component carries, for all FOVs.}
    \label{fig:rs_all_scree}
\end{figure}

\begin{figure}[H]
    \centering
    \includegraphics[width=0.72\textwidth]{summary_downloads/rs_all_pca_pathological_score.pdf}
    \caption{All FOVs, coloured by stage.}
    \label{fig:rs_all_pca}
\end{figure}

\begin{figure}[H]
    \centering
    \includegraphics[width=0.72\textwidth]{summary_downloads/rs_all_pca_organ.pdf}
    \caption{The same PCA, coloured by organ.}
    \label{fig:rs_all_pca_organ}
\end{figure}

\begin{figure}[H]
    \centering
    \includegraphics[width=0.49\textwidth]{summary_downloads/rs_all_loadings_pc1.pdf}
    \includegraphics[width=0.49\textwidth]{summary_downloads/rs_all_loadings_pc2.pdf}
    \caption{The rules behind PC1 (left) and PC2 (right).}
    \label{fig:rs_all_loadings}
\end{figure}

\begin{figure}[H]
    \centering
    \includegraphics[width=\textwidth]{summary_downloads/rs_all_fovs_pc1.pdf}
    \caption{Five FOVs along PC1, from its lowest score (A) to its highest (E), drawn as
    maps of their cells.}
    \label{fig:rs_all_fovs}
\end{figure}

\begin{figure}[H]
    \centering
    \includegraphics[width=\textwidth]{summary_downloads/rs_all_pair_fovs.pdf}
    \caption{The two FOVs whose rules are closest -- the pair squared on
    Figure~\ref{fig:rs_all_pca}. The panel above shows FOVs that should look different;
    this pair should look alike.}
    \label{fig:rs_all_pair}
\end{figure}

\begin{figure}[H]
    \centering
    \includegraphics[width=0.85\textwidth]{summary_downloads/rs_all_abundance_correlation.pdf}
    \caption{How closely each cell type's share of the cells follows PC1 and PC2. Goblet is
    the one cell type running with PC1 and the others run against it, the mark of a
    single composition gradient rather than several separate effects: Goblet is far more common in
    the colon than the duodenum, so that gradient is largely the organ.
    Every other bar stays close to the band, so no cell type dominates this PCA.}
    \label{fig:rs_all_abundance}
\end{figure}

\begin{figure}[H]
    \centering
    \includegraphics[width=0.49\textwidth]{summary_downloads/rs_all_abundance_goblet.pdf}
    \includegraphics[width=0.49\textwidth]{summary_downloads/rs_all_abundance_plasma_cd38.pdf}
    \caption{The same PCA of all FOVs, coloured by how common each of its two
    strongest cell types is: Goblet and Plasma\_CD38. Chosen from the figure above, so each part
    shows the cell types that matter in it.}
    \label{fig:rs_all_abund_cells}
\end{figure}

% ---------------------------------------------------------
\subsection{Colon, Control vs Severe}

The two extreme stages inside one organ. Neither organ nor the middle
stages can explain a split here.

\begin{figure}[H]
    \centering
    \includegraphics[width=0.62\textwidth]{summary_downloads/rs_control_severe_colon_scree.pdf}
    \caption{How much of the picture each component carries, for the colon, Control vs Severe.}
    \label{fig:rs_cs_colon_scree}
\end{figure}

\begin{figure}[H]
    \centering
    \includegraphics[width=0.62\textwidth]{summary_downloads/rs_control_severe_colon_pca_pathological_score.pdf}
    \caption{The colon FOVs scored Control or Severe. Control and Severe come apart, most
    clearly along PC2.}
    \label{fig:rs_colon_pca}
\end{figure}

\begin{figure}[H]
    \centering
    \includegraphics[width=0.49\textwidth]{summary_downloads/rs_control_severe_colon_loadings_pc1.pdf}
    \includegraphics[width=0.49\textwidth]{summary_downloads/rs_control_severe_colon_loadings_pc2.pdf}
    \caption{The rules behind PC1 (left) and PC2 (right). On PC1 the adaptive-immune rules
    sit at one end and the stromal rules at the other.}
    \label{fig:rs_colon_loadings}
\end{figure}

\begin{figure}[H]
    \centering
    \includegraphics[width=\textwidth]{summary_downloads/rs_control_severe_colon_fovs_pc1.pdf}
    \caption{Five colon FOVs along PC1, from its lowest score (A) to its highest (E).}
    \label{fig:rs_colon_fovs}
\end{figure}

The panel above walks the component end to end, so it shows FOVs that should look
\emph{different}. The pair below is the other half: the two FOVs whose rules are closest,
squared on the scatter above. They come from different patients.


\begin{figure}[H]
    \centering
    \includegraphics[width=\textwidth]{summary_downloads/rs_control_severe_colon_pair_fovs.pdf}
    \caption{The two FOVs of that pair, drawn as maps of their cells.}
    \label{fig:rs_colon_pair_fovs}
\end{figure}

\begin{figure}[H]
    \centering
    \includegraphics[width=0.85\textwidth]{summary_downloads/rs_control_severe_colon_abundance_correlation.pdf}
    \caption{How closely each cell type's share of the cells follows PC1 and PC2. Here the
    pattern is not one gradient: the immune types run strongly with PC1, matching the immune-versus-stromal split in the loadings.
    Composition and rules point the same way, so this separation cannot yet be credited to
    the rules alone. Note the wider grey band: with fewer FOVs it takes a larger rho to reach
    $p<0.05$.}
    \label{fig:rs_colon_abundance}
\end{figure}

\begin{figure}[H]
    \centering
    \includegraphics[width=0.49\textwidth]{summary_downloads/rs_control_severe_colon_abundance_cd4t.pdf}
    \includegraphics[width=0.49\textwidth]{summary_downloads/rs_control_severe_colon_abundance_mast.pdf}
    \caption{The same PCA of the colon, Control vs Severe, coloured by how common each of its two
    strongest cell types is: CD4T and Mast. Chosen from the figure above, so each part
    shows the cell types that matter in it.}
    \label{fig:rs_cs_colon_abund_cells}
\end{figure}

% ---------------------------------------------------------
\subsection{Duodenum, Control vs Severe}

The same two stages in the other organ, shown the same way as the colon
above. \textbf{Nothing separates here} -- Control and Severe stay mixed, and PC1 carries no
more variance than shuffled rules would give it, so the loadings and maps below describe an
axis that does not mean anything.

\begin{figure}[H]
    \centering
    \includegraphics[width=0.62\textwidth]{summary_downloads/rs_control_severe_duodenum_scree.pdf}
    \caption{How much of the picture each component carries, for the duodenum, Control vs Severe.}
    \label{fig:rs_cs_duo_scree}
\end{figure}

\begin{figure}[H]
    \centering
    \includegraphics[width=0.62\textwidth]{summary_downloads/rs_control_severe_duodenum_pca_pathological_score.pdf}
    \caption{The duodenum FOVs scored Control or Severe.}
    \label{fig:rs_duodenum_pca}
\end{figure}

\begin{figure}[H]
    \centering
    \includegraphics[width=0.49\textwidth]{summary_downloads/rs_control_severe_duodenum_loadings_pc1.pdf}
    \includegraphics[width=0.49\textwidth]{summary_downloads/rs_control_severe_duodenum_loadings_pc2.pdf}
    \caption{The rules behind PC1 (left) and PC2 (right).}
    \label{fig:rs_duodenum_loadings}
\end{figure}

\begin{figure}[H]
    \centering
    \includegraphics[width=\textwidth]{summary_downloads/rs_control_severe_duodenum_fovs_pc1.pdf}
    \caption{Five duodenum FOVs along PC1, from its lowest score (A) to its highest (E).}
    \label{fig:rs_duodenum_fovs}
\end{figure}

\begin{figure}[H]
    \centering
    \includegraphics[width=\textwidth]{summary_downloads/rs_control_severe_duodenum_pair_fovs.pdf}
    \caption{The two FOVs of the duodenum, Control vs Severe scope whose rules are closest -- the pair
    squared on the scatter above.}
    \label{fig:rs_cs_duo_pair}
\end{figure}

\begin{figure}[H]
    \centering
    \includegraphics[width=0.85\textwidth]{summary_downloads/rs_control_severe_duodenum_abundance_correlation.pdf}
    \caption{How closely each cell type's share of the cells follows PC1 and PC2. Every bar
    sits at or inside the band, so no cell type drives this PCA either.}
    \label{fig:rs_duodenum_abundance}
\end{figure}

\begin{figure}[H]
    \centering
    \includegraphics[width=0.49\textwidth]{summary_downloads/rs_control_severe_duodenum_abundance_cd4t.pdf}
    \includegraphics[width=0.49\textwidth]{summary_downloads/rs_control_severe_duodenum_abundance_nk.pdf}
    \caption{The same PCA of the duodenum, Control vs Severe, coloured by how common each of its two
    strongest cell types is: CD4T and NK. Chosen from the figure above, so each part
    shows the cell types that matter in it.}
    \label{fig:rs_cs_duo_abund_cells}
\end{figure}

% ---------------------------------------------------------
\subsection{Colon, Mild vs Severe}

The two disease stages without the healthy controls. Mild and Severe
stay mixed. Taken with the part above, the fingerprints tell diseased tissue from healthy but
not one degree of disease from another.

\begin{figure}[H]
    \centering
    \includegraphics[width=0.62\textwidth]{summary_downloads/rs_mild_severe_colon_scree.pdf}
    \caption{How much of the picture each component carries, for the colon, Mild vs Severe.}
    \label{fig:rs_ms_colon_scree}
\end{figure}

\begin{figure}[H]
    \centering
    \includegraphics[width=0.62\textwidth]{summary_downloads/rs_mild_severe_colon_pca_pathological_score.pdf}
    \caption{The colon FOVs scored Mild or Severe.}
    \label{fig:rs_ms_colon_pca}
\end{figure}

\begin{figure}[H]
    \centering
    \includegraphics[width=0.49\textwidth]{summary_downloads/rs_mild_severe_colon_loadings_pc1.pdf}
    \includegraphics[width=0.49\textwidth]{summary_downloads/rs_mild_severe_colon_loadings_pc2.pdf}
    \caption{The rules behind PC1 (left) and PC2 (right).}
    \label{fig:rs_ms_colon_loadings}
\end{figure}

\begin{figure}[H]
    \centering
    \includegraphics[width=\textwidth]{summary_downloads/rs_mild_severe_colon_fovs_pc1.pdf}
    \caption{Five colon FOVs along PC1, from its lowest score (A) to its highest (E).}
    \label{fig:rs_ms_colon_fovs}
\end{figure}

\begin{figure}[H]
    \centering
    \includegraphics[width=\textwidth]{summary_downloads/rs_mild_severe_colon_pair_fovs.pdf}
    \caption{The two FOVs of the colon, Mild vs Severe scope whose rules are closest -- the pair
    squared on the scatter above.}
    \label{fig:rs_ms_colon_pair}
\end{figure}

\begin{figure}[H]
    \centering
    \includegraphics[width=0.85\textwidth]{summary_downloads/rs_mild_severe_colon_abundance_correlation.pdf}
    \caption{How closely each cell type's share of the cells follows PC1 and PC2.}
    \label{fig:rs_ms_colon_abundance}
\end{figure}

\begin{figure}[H]
    \centering
    \includegraphics[width=0.49\textwidth]{summary_downloads/rs_mild_severe_colon_abundance_apc.pdf}
    \includegraphics[width=0.49\textwidth]{summary_downloads/rs_mild_severe_colon_abundance_cd4t.pdf}
    \caption{The same PCA of the colon, Mild vs Severe, coloured by how common each of its two
    strongest cell types is: APC and CD4T. Chosen from the figure above, so each part
    shows the cell types that matter in it.}
    \label{fig:rs_ms_colon_abund_cells}
\end{figure}

% ---------------------------------------------------------
\subsection{Duodenum, Mild vs Severe}

The same two stages in the other organ. \textbf{Nothing separates
here either} -- as in the duodenum part above, the two stages stay mixed and the figures
below describe an axis that does not track disease.

\begin{figure}[H]
    \centering
    \includegraphics[width=0.62\textwidth]{summary_downloads/rs_mild_severe_duodenum_scree.pdf}
    \caption{How much of the picture each component carries, for the duodenum, Mild vs Severe.}
    \label{fig:rs_ms_duo_scree}
\end{figure}

\begin{figure}[H]
    \centering
    \includegraphics[width=0.62\textwidth]{summary_downloads/rs_mild_severe_duodenum_pca_pathological_score.pdf}
    \caption{The duodenum FOVs scored Mild or Severe.}
    \label{fig:rs_ms_duodenum_pca}
\end{figure}

\begin{figure}[H]
    \centering
    \includegraphics[width=0.49\textwidth]{summary_downloads/rs_mild_severe_duodenum_loadings_pc1.pdf}
    \includegraphics[width=0.49\textwidth]{summary_downloads/rs_mild_severe_duodenum_loadings_pc2.pdf}
    \caption{The rules behind PC1 (left) and PC2 (right).}
    \label{fig:rs_ms_duodenum_loadings}
\end{figure}

\begin{figure}[H]
    \centering
    \includegraphics[width=\textwidth]{summary_downloads/rs_mild_severe_duodenum_fovs_pc1.pdf}
    \caption{Five duodenum FOVs along PC1, from its lowest score (A) to its highest (E).}
    \label{fig:rs_ms_duodenum_fovs}
\end{figure}

\begin{figure}[H]
    \centering
    \includegraphics[width=\textwidth]{summary_downloads/rs_mild_severe_duodenum_pair_fovs.pdf}
    \caption{The two FOVs of the duodenum, Mild vs Severe scope whose rules are closest -- the pair
    squared on the scatter above.}
    \label{fig:rs_ms_duo_pair}
\end{figure}

\begin{figure}[H]
    \centering
    \includegraphics[width=0.85\textwidth]{summary_downloads/rs_mild_severe_duodenum_abundance_correlation.pdf}
    \caption{How closely each cell type's share of the cells follows PC1 and PC2.}
    \label{fig:rs_ms_duodenum_abundance}
\end{figure}

\begin{figure}[H]
    \centering
    \includegraphics[width=0.49\textwidth]{summary_downloads/rs_mild_severe_duodenum_abundance_cd4t.pdf}
    \includegraphics[width=0.49\textwidth]{summary_downloads/rs_mild_severe_duodenum_abundance_endothelial.pdf}
    \caption{The same PCA of the duodenum, Mild vs Severe, coloured by how common each of its two
    strongest cell types is: CD4T and Endothelial. Chosen from the figure above, so each part
    shows the cell types that matter in it.}
    \label{fig:rs_ms_duo_abund_cells}
\end{figure}

% ---------------------------------------------------------
\subsection{Do a patient's FOVs sit together?}

Colouring a scatter by patient is not readable -- most patients have more than one FOV, and no
palette keeps that many colours apart -- so this is measured instead of drawn.
For each patient we take the mean distance between their own FOVs in the PC1--PC2 plane of
the all-FOV PCA, and divide it by the mean distance between any two FOVs; below $1$ means
that patient's FOVs sit closer together than FOVs in general.
Distances between pairs are used on purpose: unlike the distance to a patient's own centre,
the value does not drift simply because a patient has few FOVs.
Patients with a single FOV are left out, since one FOV has no spread.

\begin{figure}[H]
    \centering
    \includegraphics[width=0.85\textwidth]{summary_downloads/rs_all_patient_spread.pdf}
    \caption{One dot per patient, for the patients with two or more FOVs. Most sit
    closer together than average -- the count is in the title -- and this holds at every FOV count
    rather than only for the patients with fewest FOVs.}
    \label{fig:rs_patient}
\end{figure}

% ---------------------------------------------------------
\vspace{0.5cm}
\hrule
\vspace{0.5cm}

% ---------------------------------------------------------


\end{document}
